% ECONOMICS_PROBLEM: {"id": "H63-302", "title": "Sovereign debt limits and default", "jel_code": "H63", "source_id": "OP-0015"}
% !TeX program = xelatex
\documentclass[11pt,a4paper]{article}
\usepackage[margin=23mm]{geometry}
\usepackage{fontspec}
\setmainfont{Latin Modern Roman}
\usepackage{amsmath,amssymb,mathtools}
\usepackage{xcolor,enumitem,booktabs,longtable,array,xurl,hyperref}
\definecolor{muted}{HTML}{53636C}
\hypersetup{colorlinks=true,urlcolor=blue,linkcolor=blue}
\setlength{\parindent}{0pt}
\setlength{\parskip}{5pt}
\newcommand{\R}{\mathbb R}
\newcommand{\E}{\mathbb E}
\newcommand{\N}{\mathbb N}
\newcommand{\ind}{\mathbf 1}
\newcommand{\Var}{\operatorname{Var}}
\newcommand{\Cov}{\operatorname{Cov}}
\newcommand{\supp}{\operatorname{supp}}
\DeclareMathOperator*{\argmin}{arg\,min}
\DeclareMathOperator*{\argmax}{arg\,max}
\newcommand{\entrymeta}[3]{{\sffamily\small\color{muted}\textbf{\texttt{#1}}\quad|\quad #2\par #3\par}}
\newcommand{\Problemref}[1]{\texttt{#1}}
\newcommand{\JELref}[2]{\texttt{#2} (JEL \texttt{#1})}
\begin{document}
\section*{Sovereign debt limits and default}
\textbf{Problem ID:} \texttt{H63-302}\quad \textbf{JEL:} \texttt{H63}\par
% BEGIN ECONOMICS STATEMENT
\label{problem:OP-0015}
\label{atlas:prob:A5}

\noindent\textbf{Original atlas ID:} A5. \textbf{Formulation:} editorial specification of a research family.

\subsection*{Definitions and assumptions}

As a baseline, let income $y$ in a finite positive state space follow a specified Markov law, let one-period face-value debt $b\in[0,\bar b]$ be due now, and let the sovereign discount at $\beta\in(0,1)$. Repayment value and total value satisfy
\[
V^R(y,b)=\sup_{b'\in[0,\bar b]:\,c>0}
 \{u(c)+\beta\mathbb E[V(y',b')\mid y]\},\quad
c=y-b+q(y,b')b',\quad V=\max\{V^R,V^D\}.
\]
Set $V^R=-\infty$ if repayment is infeasible. Here $u$ is the specified utility, $V^D$ follows a specified output-loss, exclusion, and reentry law, and $q$ is the equilibrium price. With risk-neutral lenders, zero recovery, and a constant safe rate $r_f>-1$,
\[
q(y,b')=\frac{\Pr\{\text{repayment at }(y',b')\mid y\}}{1+r_f}.
\]
An equilibrium includes values, borrowing and default policies, prices, and tie-breaking at indifference.

\subsection*{Formal research question}

For a stated structural class and observed income, bond prices, default, and restructuring outcomes, characterize the identified equilibrium set. Define safe debt capacity relative to an explicit horizon $H$, tolerance $\epsilon\in[0,1)$, and initial state $z$ by
\[
b^{\rm safe}(z;H,\epsilon)=\sup\{b\in[0,\bar b]:
 \sup_{e\in\mathcal E_I}\Pr_e(\text{default by }H\mid z,b)\le\epsilon\},
\]
with a stated convention when the feasible set is empty. Determine how this capacity changes with bank exposure, creditor composition, maturity, and repayment incentives. Analyze post-default recovery using an explicitly added bargaining or restructuring model; zero recovery in the displayed baseline cannot explain recovery rates. Test any claimed monotone default threshold rather than assuming one exists.

\subsection*{Interpretation and scope}

Debt capacity is state-, horizon-, and criterion-dependent; it is not a universal debt-to-output constant. Willingness to repay and ability to service debt are distinct restrictions. Equilibrium bond pricing anticipates the future policy response, so holding prices fixed when varying borrowing generally changes the model. Domestic bank losses, rollover crises, and official-creditor bargaining require extra states, constraints, and creditor behavior. The safe-capacity set need not be an interval without monotonicity assumptions, and its supremum need not itself be safe unless appropriate closure holds. Empirical work should therefore report the feasible set or its boundaries, selection uncertainty, and assumptions translating observed spreads into default probabilities.

\subsection*{Source audit and status}

Eaton--Gersovitz (1981), Arellano (2008), and Chatterjee--Eyigungor (2012) are verified. Their models establish substantial borrowing/default and maturity results under specific restrictions; none supplies the atlas's universal safe-debt or recovery formula. The baseline Bellman equations are an editorial simplification illustrating the required primitives, rather than an exact reproduction of every cited model. The displayed lender price applies only to one-period zero-recovery debt with the stated risk neutrality. The remaining research agenda is an empirically disciplined extension integrating creditors, banks, policy incentives, and restructuring. Historical foundations do not certify that all sovereign-default benchmark characterization remains open.

\subsection*{References}

\begin{enumerate}

\item[S1] Jonathan Eaton and Mark Gersovitz (1981). \emph{Debt with Potential Repudiation: Theoretical and Empirical Analysis}. Review of Economic Studies 48(2), 289--309. \url{https://academic.oup.com/restud/article-abstract/48/2/289/1572032}. \textit{Locator:} article; endogenous sovereign borrowing and default. \textit{Establishes:} A foundational repayment-incentive model. \textit{Does not establish:} a universal safe debt-to-GDP threshold.

\item[S2] Cristina Arellano (2008). \emph{Default Risk and Income Fluctuations in Emerging Economies}. American Economic Review 98(3), 690--712. \url{https://www.aeaweb.org/articles?id=10.1257/aer.98.3.690}. \textit{Locator:} abstract; recursive default model. \textit{Establishes:} Joint determination of output, debt, and default risk in a small open economy. \textit{Does not establish:} all political, banking, and restructuring mechanisms.

\item[S3] Satyajit Chatterjee and Burcu Eyigungor (2012). \emph{Maturity, Indebtedness, and Default Risk}. American Economic Review 102(6), 2674--2699. \url{https://pubs.aeaweb.org/doi/10.1257/aer.102.6.2674}. \textit{Locator:} abstract; long-term debt model. \textit{Establishes:} Role of maturity and rollover risk in the specified environment. \textit{Does not establish:} universal debt capacity or endogenous post-default bargaining recovery.

\end{enumerate}

\subsection*{Common conventions: Source mathematical conventions}

Unless an entry states otherwise, agent and firm sets are finite. State and action spaces are measurable, strategies and outcome maps are measurable, probability laws are specified, and expectations exist for the targets under discussion. Infinite-horizon discount factors lie in $(0,1)$ and discounted objectives are finite. Norms, tolerances and complexity input sizes must be specified for computational comparisons. Constants or primitives that appear locally are inputs, not empirical facts asserted by this catalogue.

\textbf{Identification.} A statistical model $m\in\mathcal M$ implies an observable law $P_m$ and target $\theta(m)$. At law $P$, its identified set is
\[
\Theta(P;\mathcal M)=\{\theta(m):m\in\mathcal M,\ P_m=P\}.
\]
Point identification holds when this set is a singleton, or equivalently when $P_{m_1}=P_{m_2}$ implies $\theta(m_1)=\theta(m_2)$ for all compatible models. Sharp bounds mean bounds attained, or approached when the boundary is not attained, by compatible models. Writing this set is a definition, not a solution: the substantive task is to establish informative restrictions and derive or estimate the set. An unrestricted-confounding model may give only trivial bounds.

\textbf{Inference.} Identification, estimability and valid uncertainty quantification are separate questions. Fix $\alpha\in(0,1)$. A confidence procedure $C_n$ over a declared law class $\mathcal P$ has asymptotic uniform coverage at level $1-\alpha$ when
\[
\liminf_{n\to\infty}\inf_{P\in\mathcal P}\Pr_P\{\theta(P)\in C_n\}\geq1-\alpha.
\]
For set-identified targets the coverage event must be defined separately, for example coverage of the full identified set. If an entry uses identification-indexed models instead of uniquely indexed laws, the same coverage convention is applied over those models. Pointwise limits do not establish uniform validity, and asymptotic validity does not guarantee finite-sample precision. Sample sizes, dependence structures and weak-identification sequences are part of each application.

\textbf{Counterfactuals.} $Y_i(a)$ is a potential outcome under a specified intervention, including its timing, implementation and versions. It is not $Y_i$ conditional on voluntarily choosing $a$. A full intervention vector permits interference; shorthand scalar treatment notation does not silently impose no interference. Assignment, selection, attrition, anticipation and measurement must be specified. A claim about an instrument's local effect is not automatically a population policy effect.

\textbf{Economic equilibrium.} A counterfactual model includes best responses, constraints, feasibility, market clearing, state transitions and expectations consistent with the stated information. When several equilibria exist, either a declared selector is an input or the target is set-valued. Comparisons across policy regimes must use the stated selection convention. Reduced-form relationships are not presumed invariant after an equilibrium-changing intervention.

\textbf{Optimization and welfare.} Optimization is over the stated feasible set. If attainment is not established, $\sup$ or $\inf$ and $\varepsilon$-optimal choices replace an asserted maximizer or minimizer. For example, with value $V=\sup_{a\in\mathcal A}W(a)<\infty$, an $\varepsilon$-optimal set is $\{a\in\mathcal A:W(a)\geq V-\varepsilon\}$ for $\varepsilon>0$. Benefits, costs, risk and health or environmental outcomes can be aggregated only using declared units and conversion weights. Interpersonal utility weights, the treatment of fiscal transfers and foreign profits, discounting, and the behavioral welfare criterion are explicit inputs. A comparative-static derivative additionally requires existence and appropriate differentiability of the selected counterfactual.

\textbf{Sources.} Local labels [S1], [S2], and so forth restart in every question. Locators identify the relevant abstract, model, section or result at the level actually checked; a bibliographic or abstract-level verification is not represented as inspection of an entire proof. Working papers are distinguished from later publications and changing versions. Source summaries are paraphrased. All URL checks and editorial formulations refer to the audit date stated above.
% END ECONOMICS STATEMENT
\end{document}
